% !TeX root = ../sdc_regulations.tex

\color{blue}
\section[Safety Regulations]{\textcolor{blue}{Safety Regulations}}

\subsection[Pre-Match and Operational Safety]{\textcolor{blue}{Pre-Match and Operational Safety}}
\begin{itemize}
	\item{The teams are not allowed to start any motors before the official start signal is given by the judges.}
	\item{Before the match begins, each team must call "Ready" to the judges to confirm that all drones and operators are in a safe and prepared state.}
	\item{Entering the flight arena during the match is strictly prohibited for safety reasons.}
	\item{Any violation of these safety rules may result in disqualification from the current match or other disciplinary action as determined by the judges.}
	\item{Safety takes priority over scoring — teams must always operate their drones in a controlled and responsible manner.}
\end{itemize}

\subsection[Collision Prevention]{\textcolor{blue}{Collision Prevention}}
\begin{itemize}
	\item{Each drone is recommended to be equipped with a spherical protection case, which helps to protect both the drone itself and other drones from damage during the match.}
	\item{The protective spheres must be provided by the teams themselves. brigkAir does not supply drone protection equipment.}
	\item{If a drone becomes damaged or inoperative during a match, no replacement drone may be added to the team's lineup. The team must continue the match with the remaining functional drones.}
	\item{Each team may register one (1) spare drone as a backup before the tournament begins. This spare drone can only be used to replace a malfunctioning unit between matches, not during an active match.}
	\item{Teams must always avoid intentional collisions or actions that could cause damage to other drones or the arena. Fair and safe flying behavior is mandatory.}
	\item{If a team's number of operational drones drops below two, the team will be disqualified from the current match for safety and fairness reasons.}
\end{itemize}

\subsection[Continuity of C2 Link]{\textcolor{blue}{Continuity of C2 Link}}
\begin{itemize}
	\item{A continuous Command and Control (C2) link between the operator and each drone must be maintained at all times.}
	\item{If the C2 link is interrupted or lost, an automatic flight termination or controlled landing must be initiated no later than 2 seconds after the loss of connection.}
	\item{In the event of a C2 link failure, the match will be paused or safely terminated by the judges. The affected team will not be penalized, as maintaining safety is the highest priority in this competition.}
	\item{Once the drone and the area are confirmed safe, the judges will decide whether the match can be resumed or restarted, depending on the situation and remaining time.}
\end{itemize}

\subsection[Emergency Shutoff and Kill Switch]{\textcolor{blue}{Emergency Shutoff and Kill Switch}}
\begin{itemize}
	\item{Each drone must be equipped with a physical or virtual kill switch, accessible to the operator or the team's Safety Officer, and its correct functionality must be demonstrated prior to the match.}
	\item{Activation of the kill switch must result in an immediate landing or motor shutdown of the drone.}
	\item{Each team must have one master button or command that can simultaneously land or shut down all drones of the team in case of an emergency.}
	\item{Improper or unsafe operation after a kill-switch event may lead to disqualification, but safe and responsible use of the emergency system will not result in penalties.}
\end{itemize}
\color{black}
